\section*{Glossaire}
\addcontentsline{toc}{section}{Glossaire}

Les termes suivants sont utilisés dans le rapport. Ceux accompagnés d'un astérisque dans le texte sont définis ici.

\begin{description}[style=nextline]
    \item[RSE] Responsabilité Sociétale des Entreprises : démarche par laquelle l'entreprise intègre les enjeux sociaux, environnementaux et éthiques dans sa stratégie et ses opérations.
    \item[ESG] Environmental, Social and Governance : cadre d'évaluation des performances extra-financières d'une entreprise.
    \item[DEI] Diversity, Equity and Inclusion : en français, Diversité, Équité et Inclusion.
    \item[KPI] Key Performance Indicator : indicateur clé de performance permettant de suivre l'atteinte d'un objectif.
    \item[GHG] Greenhouse Gas : gaz à effet de serre (GES).
    \item[Scope 1, 2, 3] Catégories d'émissions de gaz à effet de serre définies par le GHG Protocol : Scope 1 (directes), Scope 2 (indirectes liées à l'énergie), Scope 3 (autres indirectes de la chaîne de valeur).
    \item[LCA] Life Cycle Assessment : analyse du cycle de vie, méthode d'évaluation multi-critères des impacts environnementaux d'un produit sur l'ensemble de son cycle de vie.
    \item[EDGE] Economic Dividends for Gender Equality : certification internationale en matière d'équité entre les genres en entreprise.
    \item[B4SI] Business for Societal Impact : cadre d'évaluation et d'assurance des investissements communautaires des entreprises.
    \item[SLT] Senior Leadership Team : équipe de direction exécutive d'Iveco Group.
    \item[EEHS] Energy, Environment, Health and Safety : fonction regroupant énergie, environnement, santé et sécurité.
    \item[ISO 45001] Norme internationale pour les systèmes de management de la santé et de la sécurité au travail.
    \item[CSRD] Corporate Sustainability Reporting Directive : directive européenne sur la publication d'informations de durabilité par les entreprises.
    \item[ESRS] European Sustainability Reporting Standards : normes européennes de publication d'informations de durabilité.
    \item[IT] Information Technology : technologies de l'information.
    \item[WSL] Windows Subsystem for Linux : sous-système Windows pour Linux permettant d'exécuter un environnement Linux sur Windows.
    \item[CI/CD] Continuous Integration / Continuous Deployment : intégration et déploiement continus, pratiques visant à automatiser les tests et les mises en production.
    \item[IMDS] International Material Data System : base de données sectorielle automobile répertoriant la composition des matériaux et substances des véhicules.
    \item[OS] Operating System : système d'exploitation, logiciel responsable de la gestion des ressources matérielles et logicielles d'un ordinateur.
    \item[IDE] Integrated Development Environment : environnement de développement intégré, outil regroupant éditeur, compilateur et débogueur.
    \item[PIA] Privacy Impact Assessment : analyse d'impact relative à la protection des données.
    \item[Green IT / Numérique responsable] Ensemble des pratiques visant à réduire l'empreinte environnementale et sociale des systèmes d'information.
    \item[Privacy by Design] Protection de la vie privée dès la conception : intégration de la protection des données personnelles dès les premières étapes d'un projet.
    \item[Signalement (whistleblowing)] Action de signaler, par un salarié ou un tiers, des faits répréhensibles au sein d'une organisation (anglicisme : \textit{whistleblowing}).
    \item[Bénévolat d'entreprise] Actions de volontariat menées par les collaborateurs, souvent durant leur temps de travail (en anglais : \textit{Corporate Volunteering}).
    \item[Bénévolat de compétences] Mise à disposition de savoir-faire professionnels au profit d'associations ou de communautés (en anglais : \textit{Skills-Based Volunteering}).
\end{description}
